\section{Cognitive, Emotional, and Social Foundations of Script Design}

This appendix presents academic foundations supporting the design and implementation of the Invocation of Rights, as discussed in Sections 5 and 6. Sources are organized thematically: cognitive performance under stress; emotional and behavioral factors; and social and legal behavior. Together they show that under stress people recall rather than reason, and that normalized civic behaviors enhance rights assertion.

% --- helpers ----------------------------------------------------------
\setlength{\LTpre}{3pt}
\setlength{\LTpost}{3pt}

% ----------------------------------------------------------------------
\Needspace{8\baselineskip}  % about 8 lines; adjust as desired
\subsection*{Table 1 – Cognitive performance under stress}

\begin{longtable}{@{}P{0.23\textwidth} P{0.33\textwidth} P{0.44\textwidth}@{}}
\toprule
\textbf{Author} & \textbf{Title} & \textbf{Relevance}\\
\midrule
\endfirsthead
\toprule
\textbf{Author} & \textbf{Title} & \textbf{Relevance}\\
\midrule
\endhead
\midrule
\multicolumn{3}{r}{\emph{Continued on next page}}\\
\midrule
\endfoot
\bottomrule
\endlastfoot

Arnsten (2009) &
Stress Signaling Pathways That Impair Prefrontal Cortex Structure and Function &
Details how stress impairs working memory, attention, and verbal control—faculties needed to assert rights.\\[4pt]

Sweller (1988) &
Cognitive Load During Problem Solving: Effects on Learning &
Shows that working memory collapses under complexity, validating the script’s chunked design for recall.\\[4pt]

Schwabe \& Wolf (2013) &
Stress Modulates the Engagement of Multiple Memory Systems in Classification Learning &
Explains that stress defaults to procedural memory, supporting practiced scripts over reasoning.\\[4pt]

Kahneman (2011) &
Thinking, Fast and Slow &
Distinguishes deliberate and automatic cognition, explaining why stressed civilians rely on scripted speech.\\[4pt]

Bargh \& Chartrand (1999) &
The Unbearable Automaticity of Being &
Describes how repetition builds unconscious recall, supporting automatic performance under pressure.\\[4pt]

Driskell \& Johnston (1998) &
Stress Exposure Training &
Shows that practiced verbal protocols remain accessible when higher cognition falters.\\
\end{longtable}

% ----------------------------------------------------------------------
\Needspace{8\baselineskip}  % about 8 lines; adjust as desired
\subsection*{Table 2 – Emotional and behavioral factors}

\begin{longtable}{@{}P{0.23\textwidth} P{0.33\textwidth} P{0.44\textwidth}@{}}
\toprule
\textbf{Author} & \textbf{Title} & \textbf{Relevance}\\
\midrule
\endfirsthead
\toprule
\textbf{Author} & \textbf{Title} & \textbf{Relevance}\\
\midrule
\endhead
\midrule
\multicolumn{3}{r}{\emph{Continued on next page}}\\
\midrule
\endfoot
\bottomrule
\endlastfoot

LeDoux (1996) &
The Emotional Brain: The Mysterious Underpinnings of Emotional Life &
Shows how fear impairs speech and reasoning, reinforcing the need for automatic scripts.\\[4pt]

Eysenck \textit{et al.} (2007) &
Anxiety and Cognitive Performance: Attentional Control Theory &
Explains how anxiety impairs verbal clarity, justifying simplified, rehearsed speech.\\[4pt]

Bandura (1977) &
Self‑Efficacy: Toward a Unifying Theory of Behavioral Change &
Shows that modeled behavior builds confidence, supporting civic framing of the script.\\[4pt]

Ajzen (1991) &
The Theory of Planned Behavior &
Explains how norms and perceived difficulty shape action, supporting normalization and rehearsal.\\[4pt]

Ericsson (2008) &
Deliberate Practice and the Acquisition and Maintenance of Expert Performance &
Supports rehearsal strategies, showing how structured practice builds reliable performance under pressure.\\[4pt]

Brown \& Levinson (1987) &
Politeness: Some Universals in Language Usage &
Analyzes face‑saving strategies; supports the optional “I was taught to say this” preface that mitigates confrontation.\\
\end{longtable}

% ----------------------------------------------------------------------
\Needspace{8\baselineskip}  % about 8 lines; adjust as desired
\subsection*{Table 3 – Social and legal behavior}

\begin{longtable}{@{}P{0.23\textwidth} P{0.33\textwidth} P{0.44\textwidth}@{}}
\toprule
\textbf{Author} & \textbf{Title} & \textbf{Relevance}\\
\midrule
\endfirsthead
\toprule
\textbf{Author} & \textbf{Title} & \textbf{Relevance}\\
\midrule
\endhead
\midrule
\multicolumn{3}{r}{\emph{Continued on next page}}\\
\midrule
\endfoot
\bottomrule
\endlastfoot

Rogers (2003) &
Diffusion of Innovations &
Explains how early adopters drive uptake of civic behaviors, supporting normalization of the script.\\[4pt]

Ewick \& Silbey (1998) &
The Common Place of Law: Stories from Everyday Life &
Frames rights assertion as civic behavior enacted through repeatable acts, consistent with legal consciousness research.\\[4pt]

Tyler (2006) &
Why People Obey the Law &
Shows that procedural justice improves with fair, neutral standards, supporting the script’s role in reducing escalation.\\
\end{longtable}

% ----------------------------------------------------------------------
\Needspace{8\baselineskip}  % about 8 lines; adjust as desired
\subsection*{Table 4 – Design principles reflected in the script}

\begin{longtable}{@{}P{0.25\textwidth} P{0.45\textwidth} P{0.30\textwidth}@{}}
\toprule
\textbf{Concept} & \textbf{Explanation} & \textbf{Supporting source(s)}\\
\midrule
\endfirsthead
\toprule
\textbf{Concept} & \textbf{Explanation} & \textbf{Supporting source(s)}\\
\midrule
\endhead
\midrule
\multicolumn{3}{r}{\emph{Continued on next page}}\\
\midrule
\endfoot
\bottomrule
\endlastfoot

Cognitive load &
Stress impairs working memory. The script reduces demand through chunked phrasing. &
Sweller (1988)\\[4pt]

Automaticity &
Repetition makes speech instinctive, bypassing conscious control. &
Bargh \& Chartrand (1999)\\[4pt]

Procedural memory &
Stress favors habits over reasoning. The script functions as a procedural fallback. &
Schwabe \& Wolf (2013); Driskell \& Johnston (1998)\\[4pt]

Semantic priming &
Phrases such as “invoke” echo legal language, aiding recall. &
Kahneman (2011)\\[4pt]

Approach‑avoidance conflict &
Assertiveness feels risky. Framing the script as routine lowers emotional cost. &
Bandura (1977); LeDoux (1996); Eysenck \textit{et al.}\ (2007)\\[4pt]

Normative activation &
Perceived legitimacy and cultural modeling increase use. &
Bandura (1977); Ajzen (1991); Rogers (2003)\\[4pt]

Legal consciousness &
Rights assertion becomes civic through repeatable acts. &
Ewick \& Silbey (1998)\\[4pt]

Procedural justice &
Shared scripts foster mutual understanding and reduce escalation. &
Tyler (2006)\\
\end{longtable}

% ----------------------------------------------------------------------
\Needspace{8\baselineskip}  % about 8 lines; adjust as desired
\subsection*{References}

Ajzen, I. (1991). The theory of planned behavior. \emph{Organizational Behavior and Human Decision Processes, 50}(2), 179–211. \url{https://doi.org/10.1016/0749-5978(91)90020-T}

Arnsten, A.\,F.\,T. (2009). Stress signalling pathways that impair prefrontal cortex structure and function. \emph{Nature Reviews Neuroscience, 10}(6), 410–422. \url{https://doi.org/10.1038/nrn2648}

Bandura, A. (1977). Self‑efficacy: Toward a unifying theory of behavioral change. \emph{Psychological Review, 84}(2), 191–215. \url{https://doi.org/10.1037/0033-295X.84.2.191}

Bargh, J.\,A., \& Chartrand, T.\,L. (1999). The unbearable automaticity of being. \emph{American Psychologist, 54}(7), 462–479. \url{https://doi.org/10.1037/0003-066X.54.7.462}

Brown, P., \& Levinson, S.\,C. (1987). \emph{Politeness: Some universals in language usage}. Cambridge University Press. \url{https://doi.org/10.1017/CBO9780511813085}

Driskell, J.\,E., \& Johnston, J.\,H. (1998). Stress exposure training. In J.\,A. Cannon‑Bowers \& E. Salas (Eds.), \emph{Making decisions under stress} (pp. 191–217). American Psychological Association. \url{https://doi.org/10.1037/10278-007}

Ericsson, K.\,A. (2008). Deliberate practice and the acquisition and maintenance of expert performance: A general overview. \emph{Academic Emergency Medicine, 15}(11), 988–994. \url{https://doi.org/10.1111/j.1553-2712.2008.00227.x}

Ewick, P., \& Silbey, S.\,S. (1998). \emph{The common place of law: Stories from everyday life}. University of Chicago Press. \url{https://doi.org/10.2307/2655592}

Eysenck, M.\,W., Derakshan, N., Santos, R., \& Calvo, M.\,G. (2007). Anxiety and cognitive performance: Attentional control theory. \emph{Emotion, 7}(2), 336–353. \url{https://doi.org/10.1037/1528-3542.7.2.336}

Kahneman, D. (2011). \emph{Thinking, fast and slow}. Farrar, Straus and Giroux. \url{https://doi.org/10.1007/s00362-013-0533-y}

LeDoux, J.\,E. (1996). \emph{The emotional brain: The mysterious underpinnings of emotional life}. Simon \& Schuster. \url{https://doi.org/10.1176/ajp.155.4.570}

Rogers, E.\,M. (2003). \emph{Diffusion of innovations} (5th ed.). Free Press. \url{https://doi.org/10.1016/j.jmig.2007.07.001}

Schwabe, L., \& Wolf, O.\,T. (2013). Stress modulates the engagement of multiple memory systems in classification learning. \emph{Journal of Neuroscience, 33}(26), 11042–11049. \url{https://doi.org/10.1523/JNEUROSCI.1484-12.2012}

Sweller, J. (1988). Cognitive load during problem solving: Effects on learning. \emph{Cognitive Science, 12}(2), 257–285. \url{https://doi.org/10.1207/s15516709cog1202_4}

Tyler, T.\,R. (2006). \emph{Why people obey the law}. Princeton University Press. \url{https://doi.org/10.2307/j.ctv1j66769}

\medskip\noindent
Note: These sources form the primary behavioral and social foundations for the claims made in Sections 5 and 6. A full research bibliography is available upon request.
